\selectlanguage{ngerman}%
\frenchspacing%
\begin{ledgroupsized}[r]{120mm}%
\footnotesize%
\pstart%
\noindent\textbf{Überlieferung:}%
\pend%
\end{ledgroupsized}%
\begin{ledgroupsized}[r]{114mm}%
\footnotesize%
\pstart%
\parindent -6mm%
\makebox[6mm][l]{\textit{L}}%
Aufzeichnung:
LH XXXV 10,\;16 Bl.\;8.
Ein Blatt 8\textsuperscript{o},
Ränder beschnitten;
Fragment eines \href{\Scan{LH_35_10_16_008wz}}{Wasserzeichens}:
Wiener Papier.
Eine Seite auf Bl.\;\href{\Scan{LH_35_10_16_008r}}{8\,r\textsuperscript{o}}.
Bl.\;\href{\Scan{LH_35_10_16_008v}}{8\,v\textsuperscript{o}} überliefert lediglich einen abbrechenden Briefanfang: %
\textit{Monsieur} \bracks{/} \textit{je vous dois beaucoup}
\pend%
\end{ledgroupsized}%
\vspace{5mm}%
\begin{ledgroup}%
\footnotesize%
\pstart%
\noindent%
\textbf{Datierungsgründe:}
Die Datierung der Aufzeichnung N.\;\ref{60128} beruht auf Eigenschaften des Textträgers sowie auf inhaltlichen Betrachtungen.
Siehe die ausführliche Begründung auf S.\;\pageref{60128_dat}.%
\pend%
\end{ledgroup}%
\selectlanguage{latin}%
\frenchspacing%
\vspace{8mm}%
\normalsize%
\pstart%
\noindent%
\edtext{%
\lbrack%
\href{\Scan{LH_35_10_16_008r}}%
{8\,r\textsuperscript{o}}%
\rbrack\
\edlabel{LH_35_10_16_008_008r_inveni-1}%
Inveni
qualis sit linea catenae \textit{ACB}%
\indexS{linea catenae}%
}{%
\lemma{\lbrack%
\href{\Scan{LH_35_10_16_008r}}{8~r\textsuperscript{o}}\rbrack}%
\Bfootnote{%
\textit{(1)}~Novum Genus catenariae inveni\indexS{catenaria}\indexS{genus catenariae} %
\textit{(2)}~Inveni qualis sit linea %
\textit{(a)}~catena %
\textit{(b)}~catenae \textit{ACB}%
~\textit{L}%
}}
si simpliciter catena%
\indexS{catena suspensa}
ex \textit{A} et \textit{B} suspendatur.%
\edlabel{LH_35_10_16_008_008r_inveni-2}
Ponatur jam utrinque
\edtext{%
transire catenam
super trochleis
in \textit{A} et \textit{B},%
}{%
\lemma{transire}%
\Bfootnote{%
\textit{(1)}~per %
\textit{(2)}~catenam super %
\bracks{...} et \textit{B},%
~\textit{L}%
}}
et appensa esse pondera \textit{D} et \textit{E}.%
\indexS{pondus appensum catenae}
Quia data longitudine catenae,
datur locus%
\indexS{locus centri gravitatis}
centri gravitatis,%
\indexS{centrum gravitatis catenae}
adeoque centrum gravitatis commune;%
\indexS{centrum gravitatis commune}
\edtext{%
ideo%
}{%
\lemma{ideo}%
\Bfootnote{%
\textit{erg.\;L}%
}}
definiri potest,
quousque descendere possint
\edtext{%
pondera,%
\indexS{pondus appensum catenae}
ascendere catena %
adeoque abbreviari,
ut centrum gravitatis commune catenae%
\indexS{centrum gravitatis catenae}%
\indexS{centrum gravitatis commune}
et ponderum quam maxime descendat.%
}{%
\lemma{pondera,}%
\Bfootnote{%
\textit{(1)}~abbreviari %
\textit{(2)}~ascendere %
\textit{(a)}~centrum catenae %
\textit{(b)}~catena adeoque %
\bracks{...} maxime descendat.%
~\textit{L}%
}}
\pend
\pstart
Utile erit hoc problema%
\indexS{problema utile ad mechanicen practicam}
ad mechanicen practicam%
\indexS{mechanices practica} %
\edlabel{LH_35_10_16_008_008r_anspl-1}%
ut definiatur
quanta vi opus sit %
ad catenam intendendam,%
\indexS{vis ad catenam intendendam}
ut \textit{C} ad datam altitudinem ascendat
versus horizontem \textit{AB},%
\indexS{horizon}
etsi ad eum nulla vi assignabili%
\indexS{vis assignabilis}
\edtext{%
\bracks{accedere}%
}{%
\lemma{accedere}%
\Bfootnote{%
\textit{erg. Hrsg.}%
}}
possit.%
\edlabel{LH_35_10_16_008_008r_anspl-2}
Notandum
haec etiam definire posse,
etsi pondera \textit{D} et \textit{E}%
\indexS{pondus appensum catenae}
sint
\edtext{%
inaequalia,
licet etiam \textit{A}%
}{%
\lemma{inaequalia}%
\Bfootnote{%
\textit{(1)}~. Et \textit{A} ver %
\textit{(2)}~, licet etiam \textit{A}%
~\textit{L}%
}}
et \textit{B} non sint in eodem horizonte.%
\indexS{horizon}
\pend%
\begin{edfigure}[3.00em][0.00em]
\includegraphics[width=0.380\textwidth]%
{gesamttex/edit_VIII-4/images/LH_35_10_16_008_d_008r.pdf}
\edcaption[-0.00em]{%
\bracks{\textit{Fig.\:1}}%
}%
\end{edfigure}
\newpage%
\vspace{3.0em}%
\pstart%
\noindent%
\centering%
*     *     *
\pend%
\vspace{2.0em}%
\selectlanguage{german}%
\frenchspacing%
\footnotesize%
\pstart%
\noindent%
\textls{Ausführliche Datierungsgründe:}%
\pagelabel{60128_dat}
Die titellose und undatierte Aufzeichnung N.\;\ref{60128} stellt gleichsam eine Erweiterung des auf Galilei\indexN{\textls{Galilei} (Galilaeus, Galileus), Galileo 1564\textendash1642} (\textit{Discorsi}, giornata IV: Leiden 1638, S.\;284;\cite{00050} \textit{GO} VIII, S.\;309.31--310.15)\cite{00048} zurückgehenden Problems der Kettenlinie dar:
Die gebogene waagerechte Kette, deren Linie gesucht wird, wird im Fall von N.\;\ref{60128} nicht lediglich durch ihr eigenes Gewicht belastet, sondern auch durch zwei zusätzliche, an den Enden hängende Gewichte gespannt.
Die Datierung von N.\;\ref{60128} erweist sich aus den im Folgenden angeführten Gründen als ambivalent:
\pend%
\pstart%
Im Träger der Aufzeichnung N.\;\ref{60128} liegt das Fragment eines \href{\Scan{LH_35_10_16_008wz}}{Wasserzeichens} vor, das sich eindeutig Wiener\indexO{Wien} Papier zuordnen lässt;
Papier dieser Art ist im Leibniz-Nachlass insbesondere (aber möglicherweise nicht ausschließlich) für den Zeitraum von Leibnizens letztem Aufenthalt in der habsburgischen Hauptstadt\indexO{Wien} belegt, d.\,h. von Mitte Dezemeber 1712 bis zum 3.\;September 1714 (\textit{Chronik}, S.\;233--247).\cite{01236}
Dies ist Grund für die Annahme, dass auch die Aufzeichnung N.\;\ref{60128} etwa zur gleichen Zeit entstanden sein dürfte, d.\,h. wohl nicht vor Ende 1712 und nicht nach September 1714.
Diese Vermutung wird dadurch bekräftigt, dass das auf Bl.\;\href{\Scan{LH_35_10_16_008v}}{8\,v\textsuperscript{o}} überlieferte Satzfragment (\textit{Monsieur} \bracks{/} \textit{je vous dois beaucoup}) dem Anfang des Briefes ähnelt, den Leibniz am 25.\;April 1713 aus Wien an Johann Theodor von Imbsen\indexN{\textls{Imbsen} (Imsen, Immesen), Johann Theodor von 1678\textendash1742} sandte: \textit{Je vous dois beaucoup de remercimens des temoignages effectifs de votre bonté ...} (LH XLI\;9 Bl.\;39\,v\textsuperscript{o}).\cite{02537}
\pend%
\pstart%
Es gilt jedoch zu beachten, dass mindestens ein älterer Brief Leibnizens mit ähnlichem Anfang überliefert ist, nämlich sein Schreiben vom 18./28.\;November 1692 an Christoph von Weselow:\indexN{\textls{Weselow} (Weselau), Christoph von, gest. 1695} \textit{Je vous dois beaucoup de remercimens de la bonté que vous avés eue ...} (\textit{LSB} I,\;8 N.\;333, S.\;543.5).\cite{02538}
Ferner ist nicht auszuschließen, dass Wiener\indexO{Wien} Papier der gleichen Art auch früher in Verwendung gewesen sein kann;
diese Möglichkeit wäre selbst dann nicht gebannt, wenn eindeutige Belege dafür im Leibniz-Nachlass fehlen sollten.
Folglich ist auch nicht auszuschließen, dass die Aufzeichnung N.\;\ref{60128} bereits während eines vorhergehenden Aufenthaltes Leibnizens in Wien\indexO{Wien} entstanden sein könnte, nämlich %
vom 8.\;Mai 1688 bis etwa zum 12.\;Februar 1689,
von Ende April bis Mitte Mai 1690,
vom 29.\;Oktober 1700 bis spätestens zum 20.\;Dezember 1700
oder
im Dezember 1708 (\textit{Chronik}, S.\;89--95, 103, 167--169 u.\;211\,f.).\cite{01236}
\pend%
\pstart%
Für eine frühere Entstehung spricht etwa, dass gegen Ende der Aufzeichnung N.\;\ref{60128} (\refpassage{LH_35_10_16_008_008r_anspl-1}{LH_35_10_16_008_008r_anspl-2}) derselbe Gedanke aufgegriffen wird wie in der auf die Zeitspanne zwischen Herbst 1679 und 1682 (unter Vorbehalt) datierbaren, in diesem Band veröffentlichten Aufzeichnung N.\;\ref{60198} \textit{Nulla chorda pondere praedita exacte tendi potest} (siehe die ausführlichen Datierungsgründe auf S.\;\pageref{60198_dat}).
\pend%
\pstart%
Als besonders relevant erweist sich aber, dass Leibniz im ersten Satz der Aufzeichnung N.\;\ref{60128} zwar darauf anspielt, das Problem der Kettenlinie gelöst zu haben (\textit{Inveni qualis sit linea catenae ACB si simpliciter catena ex A et B suspendatur}, \refpassage{LH_35_10_16_008_008r_inveni-1}{LH_35_10_16_008_008r_inveni-2}),
mit keinem Wort aber seine diesbezüglichen Veröffentlichungen erwähnt, auch nicht seine erste: \textit{De linea in quam flexile se pondere proprio curvat} (\textit{AE}, Juni 1691, S.\;277--281; \textit{LMG} VI, S.\;243--247).\cite{02033}\cite{01043}
Dieser Umstand kann als weiterer Beleg für eine frühere Entstehung der Aufzeichnung N.\;\ref{60128} gedeutet werden.
Denn das Problem der Kettenlinie war von Johann Bernoulli\indexN{\textls{Bernoulli}, Johann 1667\textendash1748} im Maiheft 1690 der \textit{Acta eruditorum} als Herausforderung für die mathematisch versierte Leserschaft der Zeitschrift formuliert worden:
\textit{Invenire quam curvam referat funis laxus et inter duo puncta fixa libere suspensus} %
(\textit{AE} 1690, S.\;219)\cite{02620}\cite{01023}.
In der ersten Hälfte Mai 1690 befand sich Leibniz in Wien,\indexO{Wien} auf dem Rückweg aus Italien\indexO{Italien} nach Hannover.\indexO{Hannover}
Im folgenden Juliheft der \textit{Acta eruditorum} veröffentlichte er trotzdem eine Antwort auf Bernoullis Herausforderung, in welcher er ankündigte, zwar die Lösung gefunden zu haben, sie aber erst ein Jahr später öffentlich bekanntmachen zu wollen (\textit{AE} 1690, S.\;360).\cite{02621}\cite{01023}
Der Brief an Rudolf Christian von Bodenhausen\indexN{\textls{Bodenhausen} (Bodenus, Bodenausen), Rudolf Christian, Freiherr v. 1698} vom 6.\;Juli 1690 liefert einen Überblick über Leibnizens Beschäftigung mit dem Problem der Kettenlinie und zugleich einen Beleg dafür, dass seine Lösung Anfang Juli 1690 bereits vorlag: \textit{Die linea dependirt a quadratura Hyperbolae, ist also solvirt gnug, et construiret per logarithmos} (\textit{LSB} III,\;4 N.\;264, S.\;527.5--21, bes. Z.\;12\,f.).\cite{02622}
\pend%
\pstart%
Leibniz hat offenbar noch vor seiner Abfahrt aus Wien\indexO{Wien} um Mitte Mai 1690 oder kurz danach Gelegenheit gehabt, das Maiheft der \textit{Acta eruditorum} zu lesen und sich Johann Bernoullis\indexN{\textls{Bernoulli}, Johann 1667\textendash1748} Fragestellung zu widmen.
Die Vermutung liegt somit nahe, dass er zu dieser Zeit auch die Aufzeichnung N.\;\ref{60128} verfasst hat, deren erster Satz (\refpassage{LH_35_10_16_008_008r_inveni-1}{LH_35_10_16_008_008r_inveni-2}) eine Aspielung auf die gerade gefundene, aber noch nicht veröffentlichte Lösung des Problems der Kettenlinie sein könnte.
Dies würde Leibnizens Schweigsamkeit über seinen erst ein Jahr später erschienenen Aufsatz \textit{De linea in quam flexile se pondere proprio curvat}\cite{02033} und über die weiteren Veröffentlichungen zu diesem Thema erklären.
Auch die Wiener\indexO{Wien} Herkunft des Trägers von N.\;\ref{60128} würde sich in diesem Fall als plausibel erweisen:
Wenn die Aufzeichnung in den Wochen zwischen Mitte Mai und Ende Juni 1690 entstand, als Leibniz sich noch in der habsburgischen Hauptstadt\indexO{Wien} aufhielt oder diese gerade verlassen hatte, dann ist nicht verwunderlich, dass er noch Wiener\indexO{Wien} Papier dafür verwendet hat.
\pend%
\pstart%
Angesichts der soeben dargelegten alternativen Möglichkeiten bietet sich an, für die Aufzeichnung N.\;\ref{60128} eine zweifache Datierung vorzuschlagen: \textit{Mitte Mai bis Ende Juni 1690} oder \textit{Ende 1712 bis September 1714} (beide natürlich unter Vorbehalt).
Für die spätere Datierung plädieren die bekannten Belege des Wassezeichens;
für die frühere der dargestellte inhaltliche Sachverhalt.
\pend%
\selectlanguage{latin}%
\frenchspacing%
\normalsize%
